\clearpage
\VerdierPhysicalPageAnchor{0054}%
\VerdierSemanticAnchor{ch1-1-5-4-eq1-condition}
は、$\lambda,\mu\in\mathbb Z^{n-1}$ または
$\lambda,\mu\in\mathbb Z e_n$ のとき $\alpha(\lambda,\mu)$-可換である。
$\sigma\in\mathbb Z^{n-1}$、$\tau\in\mathbb Z e_n$ に対し、
$Q(\sigma,\tau)$ で次の性質を表す：

\VerdierSemanticAnchor{ch1-1-5-4-property-q}
\begin{quote}
  \noindent\textup{$(Q(\sigma,\tau))$}\quad
  \emph{次の図式は
  $\alpha(\sigma,\tau)\alpha^{-1}(\tau,\sigma)$-可換である：}
\end{quote}

\VerdierSemanticAnchor{ch1-1-5-4-eq2}
\begin{equation}
  \resizebox{0.98\textwidth}{!}{%
    $\begin{tikzcd}[ampersand replacement=\&,column sep=4.5em,row sep=3.2em]
      FT(\sigma+\tau)
        \arrow[r,"{F*c(\sigma,\tau)}"]
        \arrow[d,"{F*c(\tau,\sigma)}"]
        \& FT(\tau)T(\sigma)
        \arrow[r,"{p(\tau)*T(\sigma)}"]
        \& T'(\tau)FT(\sigma)
        \arrow[d,"{T'(\tau)*p(\sigma)}"'] \\
      FT(\sigma)T(\tau)
        \arrow[d,"{p(\sigma)*T(\tau)}"]
        \& {}
        \& T'(\tau)T'(\sigma)F
        \arrow[d,"{c'(\sigma,\tau)^{-1}*F}"'] \\
      T'(\sigma)FT(\tau)
        \arrow[r,"{T'(\sigma)*p(\tau)}"']
        \& T'(\sigma)T'(\tau)F
        \arrow[r,"{c'(\tau,\sigma)^{-1}*F}"']
        \& T'(\sigma+\tau)F
    \end{tikzcd}$%
  }\;.
  \tag*{\(\textup{(1.5.4.2)}_{\sigma,\tau}\)}
\end{equation}

仮定により、性質 $Q(e_i,e_n)$、$Q(e_i,0)$（$1\leq i\leq n-1$）、
$Q(0,e_n)$ は成り立つ。さらに
$\alpha(\lambda,\mu)$-可換図式
$\textup{(1.5.4.1)}_{\lambda,\mu}$ を用いると、(1.5.4.3) が容易に
確認できる：

\VerdierSemanticAnchor{ch1-1-5-4-eq3}
\begin{quote}
  \noindent\textup{(1.5.4.3)}\quad
  \emph{$\sigma\in\mathbb Z^{n-1}$、
  $\tau,\tau'\in\mathbb Z e_n$ とする。性質
  $Q(\sigma,\tau+\tau')$、$Q(\sigma,\tau)$、$Q(\sigma,\tau')$
  のうち二つが成り立てば、残る一つも成り立つ。同様に、
  $\sigma,\sigma'\in\mathbb Z^{n-1}$、$\tau\in\mathbb Z e_n$ とする。
  性質 $Q(\sigma+\sigma',\tau)$、$Q(\sigma,\tau)$、
  $Q(\sigma',\tau)$ のうち二つが成り立てば、残る一つも成り立つ。}
\end{quote}

\VerdierSemanticAnchor{ch1-1-5-4-q-conclusion}
したがって性質 $Q(\sigma,\tau)$ は、すべての
$\sigma\in\mathbb Z^{n-1}$ と $\tau\in\mathbb Z e_n$ に対して成り立つ。

\VerdierSemanticAnchor{ch1-1-5-4-p-omega}
各 $\omega\in\mathbb Z^n$ を
$\omega=\sigma+\tau$（$\sigma\in\mathbb Z^{n-1}$、
$\tau\in\mathbb Z e_n$）と書く。同型
$p(\omega)\colon FT(\omega)\longrightarrow T'(\omega)F$ を、図式
$\textup{(1.5.4.1)}_{\sigma,\tau}$ を $\alpha(\sigma,\tau)$-可換にする
唯一の同型として定める。これにより、任意の対
$\lambda,\mu\in\mathbb Z^n$ について図式
$\textup{(1.5.4.1)}_{\lambda,\mu}$ を書き、次の性質を考えられる：

\VerdierSemanticAnchor{ch1-1-5-4-property-p}
\begin{quote}
  \noindent\textup{$(P(\lambda,\mu))$}\quad
  \emph{図式 $\textup{(1.5.4.1)}_{\lambda,\mu}$ は
  $\alpha(\lambda,\mu)$-可換である。}
\end{quote}
